% !TeX root = ../memoire.tex

\chapter*{Résumé}
\addcontentsline{toc}{chapter}{Résumé}

À Montpellier, les archives de l'enseignement médical sont demeurées au plus près de l'institution qui les a produites et transmises. Cette permanence recouvre pourtant une histoire mouvementée : pendant huit siècles, les documents ont circulé entre officiers et professeurs, quitté les coffres de la Faculté pour les besoins de son administration et de ses procès, disparu dans des successions, puis été retrouvés, rassemblés et reclassés. Le fonds conservé aujourd'hui ne constitue donc pas seulement la mémoire de l'institution ; son organisation porte la trace de ses transformations.

Cette conservation historique doit désormais s'inscrire dans le cadre juridique des archives publiques. Alors que l'Université de Montpellier prépare une demande de dérogation à l'obligation de versement aux Archives départementales de l'Hérault, la campagne engagée depuis 2023 vise à garantir la connaissance, la conservation, la communication et la valorisation du fonds. Réalisé au sein de la Mission Archives historiques de la Faculté de médecine, ce mémoire accompagne cette démarche en interrogeant les conditions de sa diffusion numérique. L'étude des institutions et des classements successifs précède ainsi un diagnostic des fonds, des instruments de recherche et des dispositifs existants, complété par des expérimentations et retours d'expérience professionnels.

Parti de l'hypothèse d'un portail propre aux Archives historiques, le travail aboutit à une conclusion différente : conserver à demeure ne signifie pas isoler, pas plus que diffuser ne suppose de tout réunir dans un même outil. L'accès repose sur une architecture documentaire distribuée, dans laquelle descriptions, reproductions et contenus éditoriaux demeurent confiés à des environnements distincts mais reliés par des standards, des identifiants et des responsabilités établies. L'intégration des archives au patrimoine universitaire et aux réseaux de la recherche dépend dès lors moins de la concentration des ressources que de la circulation de leurs données, de leur mise en relation avec d'autres collections et de la coopération entre les institutions qui les conservent, les décrivent et les étudient.

\vspace{0.6cm}

\noindent\textbf{Mots-clés :} Faculté de médecine de Montpellier ; archives universitaires ; archives publiques ; histoire des archives ; dérogation au versement ; description archivistique ; diffusion numérique ; patrimoine universitaire ; interopérabilité.

\vspace{1cm}

\noindent\textbf{Abstract}

In Montpellier, the archives of medical education have remained closely connected to the institution that produced and transmitted them. This apparent continuity nevertheless conceals a complex history. Over eight centuries, records circulated between officers and professors, left the Faculty's custody for administrative and legal purposes, disappeared into private estates, and were subsequently recovered, gathered and repeatedly rearranged. The present fonds is therefore more than the institutional memory of the Faculty: its organisation bears the traces of the successive transformations through which it has been preserved.

This historical custody must now be brought into line with the legal framework governing French public archives. As the University of Montpellier prepares an application for exemption from the statutory transfer of its permanent records to the Hérault Departmental Archives, the programme launched in 2023 seeks to secure the long-term preservation, accessibility and promotion of the fonds. Conducted within the Historical Archives Mission of the Faculty of Medicine, this dissertation contributes to that process by examining the conditions of digital dissemination. It combines a historical study of the successive institutions and archival arrangements with an assessment of the fonds, finding aids and existing services, supported by technical experiments and professional feedback.

The study began with the prospect of a dedicated portal but reached a different conclusion: preserving archives in their historical setting does not mean isolating them, just as disseminating them does not require gathering every resource in a single system. Access depends on a distributed documentary architecture in which descriptions, digital reproductions and editorial content remain managed in distinct environments, connected through standards, identifiers and clearly assigned responsibilities. The integration of the archives into university heritage and research networks therefore relies less on concentrating resources than on circulating their data, connecting them with other collections, and sustaining cooperation between the institutions that preserve, describe and study them.

\vspace{0.6cm}

\noindent\textbf{Keywords:} Montpellier Faculty of Medicine; university archives; public archives; archival history; exemption from archival transfer; archival description; digital dissemination; university heritage; interoperability.

\vfill

\noindent\textbf{Informations bibliographiques :}

\noindent MARTIN Clara, \emph{À demeure, en partage : conservation in situ et diffusion en réseau des archives historiques d'une institution universitaire. Le cas de la Faculté de médecine de l'Université de Montpellier}, mémoire de master « Technologies numériques appliquées à l'histoire », sous la direction d'Emmanuelle Bermès, École nationale des chartes, 2026.